\textbf{Purpose:} To test whether detailed crowd-vote distributions improve selective facial-expression prediction beyond simpler target softening, and diagnose reliability when prediction and annotation vocabularies differ. \textbf{Methods:} Twelve channel-gated Swin-T runs compare hard, tie-aware hard, dominant-confidence-preserving uniform-minority, and full soft targets across three matched seeds. Exact-pixel isolation separates training, checkpoint selection, calibration, and testing. Selected and fixed-epoch checkpoints are evaluated. Complete-vote disagreement is partitioned into unsupported mass, supported ambiguity, and excess error; calibration half-splits and a conservative fixed-grid comparator examine threshold transfer. \textbf{Results:} At 80\% coverage on 3,275 test images, mean disagreement is 26.30\% (hard), 25.82\% (tie-aware), 24.06\% (uniform-minority), and 23.68\% (soft). Soft minus uniform-minority is -0.38 percentage points (conditional 95\% pixel-cluster interval [-0.64, -0.07]). The test annotation floor exceeds calibration's by 2.97 percentage points, including 2.15 points of unsupported mass. \textbf{Conclusion:} Full soft targets yielded lower selective disagreement than the ambiguity-preserving uniform-minority control under this matched protocol. Confidence interpretation depends on the vote vocabulary and retained population. Findings concern one reused image source, not independent-domain reliability or communication benefit.
